\subsection{紧形式的共轭性}

定理 1. 设 \(\mathfrak{a}\) 是复半单李代数。

\begin{enumerate}
\def\labelenumi{\alph{enumi})}
\item
  \(\mathfrak{a}\) 有紧（相应地，可分裂）实形式。
\item
  群 \(\operatorname{Int}(\mathfrak{a})\) 传递地作用在 a 的紧（相应地，可分裂）实形式所成的集上。
\end{enumerate}

设 \(\mathfrak{h}\) 是 \(\mathfrak{a}\) 的一个嘉当子代数。则 \((\mathfrak{a},\mathfrak{h})\) 是分裂的（第 VIII 章，§2，第 1 节，注 2），并且有一个谢瓦莱系 \((X_{\alpha})\) （第 VIII 章，§4，第 4 节，命题 5 的推论）。部分 \(a)\) 于是由命题 2 得出。设 \(\mathfrak{g}\) 是 \(\mathfrak{a}\) 的紧实形式；我们证明存在 \(v\in\mathrm{Int}(\mathfrak{a})\) 使得 \(v(\mathfrak{a}_{u})=\mathfrak{g}\) 。设 \(\mathfrak{t}\) 是 \(\mathfrak{g}\) 的一个嘉当子代数；则 \(\mathfrak{t}_{(\mathbf{C})}\) 是 \(\mathfrak{a}\) 的嘉当子代数；由于 \(\mathrm{Int}(\mathfrak{a})\) 传递地作用在 \(\mathfrak{a}\) 的嘉当子代数所成的集上（第 VII 章，§3，第 2 节，定理 1），我们可化到 \(\mathfrak{t}_{(\mathbf{C})}=\mathfrak{h}\) 的情形。由于 \(\mathfrak{g}\) 是紧形式，自同态 ad \(h\) （\(h\in\mathfrak{t}\) ）的特征值都是纯虚的（ \(\S1\) ，第 3 节，命题 1），故根 \(\alpha\in\mathbb{R}\) 把 \(\mathfrak{t}\) 映到 \(i\mathbf{R}\) 中；这就蕴含 \(\mathfrak{t}=i\mathfrak{h}_{0}\) 。于是由第 2 节命题 3，存在 \(v\in\mathrm{Int}(\mathfrak{a})\) 使得 \(v(\mathfrak{a}_{u})=\mathfrak{g}\) ，从而在紧形式的情形得到 \(b)\)。最后，设 \(\mathfrak{m}_{1}\) 与 \(\mathfrak{m}_{2}\) 是 \(\mathfrak{a}\) 的两个可分裂实形式。存在标架 \((\mathfrak{m}_{1},\mathfrak{h}_{1},\mathrm{B}_{1},(X_{\alpha}^{1}))\) 与 \((\mathfrak{m}_{2},\mathfrak{h}_{2},\mathrm{B}_{2},(X_{\alpha}^{2}))\) （第 VIII 章，§4，第 1 节）。它们以显然的方式扩张成基 \(e_{1}\) 与 \(e_{2}\) （\(\mathfrak{a}\) 的）。\(\mathfrak{a}\) 的把 \(e_{1}\) 映为 \(e_{2}\) 的自同构把 \(\mathfrak{m}_{1}\) 映为 \(\mathfrak{m}_{2}\) ；于是，只需应用第 VIII 章 §5 第 3 节命题 5，便得到存在元素 \(u\) （属于 \(\mathrm{Aut}_{0}(\mathfrak{a})=\mathrm{Int}(\mathfrak{a})\) ）使得 \(u(\mathfrak{m}_{1})=\mathfrak{m}_{2}\) 。

注. 我们将在很后面给出复半单李代数的实形式的一般分类。

推论 1. 设 \(\mathfrak{g}\) 与 \(\mathfrak{g}'\) 是两个紧实李代数。则 \(\mathfrak{g}\) 与 \(\mathfrak{g}'\) 同构当且仅当复李代数 \(\mathfrak{g}_{(\mathbf{C})}\) 与 \(\mathfrak{g}'_{(\mathbf{C})}\) 同构。

条件显然是必要的。反之，假定 \(\mathfrak{g}_{(\mathbf{C})}\) 与 \(\mathfrak{g}'_{(\mathbf{C})}\) 同构。设 c（相应地，\(c'\) ）是 g（相应地，\(g'\) ）的中心，s（相应地，\(s'\) ）是 g（相应地，\(g'\) ）的导代数。则 \(\mathfrak{c}_{(\mathbf{C})}\) 与 \(\mathfrak{c}'_{(\mathbf{C})}\) 分别是 \(\mathfrak{g}_{(\mathbf{C})}\) 与 \(\mathfrak{g}'_{(\mathbf{C})}\) 的中心，因而同构；由此得出交换代数 c 与 \(c'\) 同构。同样，\(\mathfrak{s}_{(\mathbf{C})}\) 与 \(\mathfrak{s}'_{(\mathbf{C})}\) 同构，故 s 与 \(s'\) ——它们是两个同构的复半单李代数的紧实形式——由定理 1 b) 得知同构。

推论 2. 设 \(\mathfrak{a}\) 是复李代数。下列条件等价：

\begin{enumerate}
\def\labelenumi{(\roman{enumi})}
\item
  \(\mathfrak{a}\) 是约化的。
\item
  存在紧实李代数 \(\mathfrak{g}\) 使得 \(\mathfrak{a}\) 同构于 \(\mathfrak{g}(\mathbf{C})\) 。
\item
  存在紧李群 \(\mathbf{G}\) 使得 \(\mathfrak{a}\) 同构于 \(\mathrm{L}(\mathrm{G})_{(\mathbf{C})}\) 。
\end{enumerate}

由 §1 第 3 节定义 1，条件 (ii) 与 (iii) 等价且蕴含 (i)。若 \(\mathfrak{a}\) 约化，则它是一个交换代数（它显然有紧实形式）与一个半单代数（由定理 1 \(a\) ) 它有紧实形式）的直积，故 (i) 蕴含 (ii)。

推论 3. 设 \(\mathfrak{a}_1\) 与 \(\mathfrak{a}_2\) 是两个复半单李代数。\(\mathfrak{a}_1 \times \mathfrak{a}_2\) 的紧实形式是诸积 \(\mathfrak{g}_1 \times \mathfrak{g}_2\) ，其中对 \(i = 1, 2\) ，\(\mathfrak{g}_i\) 是 \(\mathfrak{a}_i\) 的紧实形式。

事实上，存在紧实形式 \(g_{1}\) （相应地，\(g_{2}\) ）——\(a_{1}\) （相应地，\(a_{2}\) ）的——；于是 \(g_{1} \times g_{2}\) 是 \(a_{1} \times a_{2}\) 的紧实形式。把定理 1 b) 应用于 \(a_{1}, a_{2}\) 与 \(a_{1} \times a_{2}\) ，即得本推论。

注意，由上面的推论 3 可知，紧实李代数 g 单当且仅当复李代数 \(\mathfrak{g}_{(\mathbf{C})}\) 单。我们说 g 是 \(A_{n}\) 型、或 \(B_{n}, \ldots\) 型的，如果 \(\mathfrak{g}_{(\mathbf{C})}\) 是 \(A_{n}\) 型、或 \(B_{n}, \ldots\) 型的（第 VIII 章，§2，第 2 节）。由上面的推论 1，两个紧单实李代数同构当且仅当它们同型。

设 G 是概单连通紧李群（第 III 章，§9，第 8 节，定义 3）。我们说 G 是 \(A_{n}\) 型、或 \(B_{n}\) 型、\ldots，如果 G 的李代数是 \(A_{n}\) 型、或 \(B_{n}\) 型、\ldots。两个单连通概单紧李群同构当且仅当它们同型。

